
% To be included
\chapter{误差的控制}

在上一章里，我们已经介绍了自行镇定的原理：由于有一个简单的装置，每当系统处于不稳定的时候，这个装置就会改变系统的性能，使系统能够不受那些偶然的误差以及元件发生无法预料的失效现象的影响。因为一个自行镇定的系统能够自动地达到稳定状态，所以，设计这种控制系统的时候，实际上，那些不稳定的运动状态场和那些稳定的运动状态场同样地被引进系统中来了。换句话说，在设计一个自行镇定的系统的过程中，我们并不企图把稳定性和不稳定性区分开来，也不设法把错误的运动状态场与正确的运动状态场加以区分。在那里，我们只是从概率的观点处理运动状态的误差，并没有采用其他的处理方法。本章中，我们将从另外一种不同的观点讨论一个复杂系统的可靠性问题。我们将要故意把误差引入系统中，然后提出问题：系统应该如何设计才能使系统可以不顾误差的影响而还可以给出合适的性能。那就是，我们希望知道如何控制误差。

关于控制误差的问题现在还处于发展的初期。目前还只能对那些最简单的作用关系讨论误差的控制，现在的理论全都是由冯·诺伊曼(J. von Neumann)\footnote{J. von Neumann, ``Probabilistic Logics and the Synthesis of Reliable Organisms from Unreliable Components'',这是作者在美国加州理工学院(California Institute of Technology)作过的一篇讲演的印本。出版于 1952 年，现已收入下列文集中:\textit{Automata Studies}, Princeton University Press, Princeton, N. J. (1956)，这个文集已有俄文译本[60]。}提出来的。我们这一章的讨论是冯·诺伊曼工作的一个说明。这样做的目的在于引入这个非常重要的题目，并且指出需要在这方面进行更进一步的研究。

\section{用加倍的办法改进可靠性}

根据一般的常识，用简单的加倍办法常常可以增加系统的可靠程度。例如，以图\ref{fig:18_1}(a)所表示的系统而论，假设它有这种特性：如果它的运转失效仅仅是没有输出而已，这时为了保证系统有较小的失效概率，我们用加倍的办法可以把同样的系统并联起来，如图\ref{fig:18_1}(b)。假设原来系统失效的概率等于 $p$（$p$ 是一个在零与 $1$ 之间的数目），于是在并联的系统中每一个单元失效的概率也是 $p$。如果并联而成的系统中各个单元相互无关，只有当每一个单元都失效的时候并联而成的系统才会失效。并联而成的系统失效的概率于是就等于 $p^n$。使倍数 $n$ 增加，我们就可以使得这个概率非常小。

\begin{figure}[h!]
\centering
\includegraphics[width=0.5\textwidth]{图18.1.png} 
\caption{}
\label{fig:18_1}
\end{figure}

但是，一般说来，控制系统中的某一个元件的失效，并不使它的输出等于零。与上一段的情形相比较，这时的失效影响会更坏些，这时系统仍然还有一个输出，但是这是个错误的输出，因此，像上面谈到的加倍方法，就不可能解决问题，因为一个错误的信号混杂在正确的信号中其结果仍然是错误信号。因此，对于防止失效说来，如果系统中有这种失灵情况，并联而成的系统失灵的概率和单独一个系统失灵的概率相等；结果，可靠性并没有得到改善。所以误差的控制问题并不像最初想像那么简单而是一个比较深刻、比较困难的问题。虽然如此，以后还会看到，加倍的原理，也即需要增加元件数目这一原理仍然是基本的想法。主要的问题在于如何把这些元件按新的方式组合起来才能有效地控制误差，因为像\ref{fig:18_1}(b)那样简单的并联组合并不总是有效的。

\section{基本元件}

为了分析起来简单起见，我们不来讨论输入和输出都是连续变数的情形，我们选择一个简单的元件作为对象，它的输出和输入只可能取两个离散值：$0$ 或者 $1$，也就是说，输入或者是“开”（激发）或者是“关”，输出或者是“开”（激发）或者是“关”。这时元件的特性由输入和输出的关系所决定。此后我们总是假定元件只能有一个输出，但是可以有几个输入。这里输出和那些输入之间也可以有一个时滞，这里，时滞的意义是：输入激发后，一定要经过一定的时间输出才激发。所以，这样一种元件具有继电器线路的性质。

\begin{wrapfigure}{r}{0.25\textwidth}
    \centering
    \includegraphics[width=0.25\textwidth]{图18.2.png}
    \caption{ }
    \label{fig:18_2}
\end{wrapfigure}
为了描述元件的特征，我们引入四种类型的输入：刺激输入；抑制输入；恒定刺激输入；以及一类永不起作用的输入，也就是接地输入。它们是用图\ref{fig:18_2}的符号表示的。这里，用一个圆圈表示元件，输入画在圆圈的左端，输出画在圆圈的右端。而且在圆圈中间有一个数字 $k$，$k$ 的意义就是：起刺激作用的输入的数目必须至少比起抑制作用的输入的数目多 $k-1$ 个，系统的输出才是开。图\ref{fig:18_3}(a)所表示元件，只有当输入 $a$ 和 $b$ 都是开的时候，输出才是开，这可以叫做 $ab$ 元件。图\ref{fig:18_3}(b)表示另一个元件，如果输入 $a$ 和 $b$ 之中有一个是开那么它的输出是开，这可以叫做 $a+b$ 元件。图\ref{fig:18_3}(c)表示只有当输入是关的时候输出才是开的元件，可以叫做 $a^{-1}$ 元件。
\begin{figure}[h!]
\centering
\includegraphics[width=0.7\textwidth]{图18.3.png} 
\caption{}
\label{fig:18_3}
\end{figure}
如果我们考虑的输入是表示某个断言是正确（开）还是不正确（关）的条件，那么，图\ref{fig:18_3}的三个元件恰好表示布尔(Boole)代数的三个基本运算（这些元件有时候也叫做基本逻辑元件）。在一个采用二进位制的数字计算机中，这些元素是运算的基本元件。如果计算运算中需要记忆装置，这种装置可以由图\ref{fig:18_3}的元件经过反馈而得到，如图\ref{fig:18_4}那样。原因是这样的：输入 $a$ 激发了一次以后，不论在下一次运算中 $a$ 是开还是关，输入总是开的。

\begin{wrapfigure}{r}{0.25\textwidth}
    \centering
    \includegraphics[width=0.25\textwidth]{图18.4.png}
    \caption{ }
    \label{fig:18_4}
\end{wrapfigure}
对于以后的讨论来说，同时用三种基本元件去处理问题有些不方便。但是，所有这三种元件实际上是一种基本元件的特殊情形。考虑图\ref{fig:18_5}所表示的\underline{谢弗元件}(Scheffer stroke)。因为两条恒定刺激输入线是一起起作用，我们可以把他们从图中去掉，把谢弗元件用下面的图表示。输出在这样的情形下是开，如果 $a$ 和 $b$ 都不是开，或者 $a,b$ 中有一个是开，输出就是开的。但是如果 $a$ 和 $b$ 都开那么输出就是关。上一段谈到的三个基本元件可以用谢弗元件按图\ref{fig:18_6}的方式构成。$ab$ 元件和 $a+b$ 元件都包含两个串联的谢弗元件；所以这两种元件的时滞自然等于那种只包含一个谢弗元件的 $a^{-1}$ 元件的时滞的两倍。但是因为我们对于时滞的了解只在于指出输入的作用先于输出若干时间，时滞真正的大小并不重要。因此，我们把只在时滞的大小上有差别的运算考虑作等价的运算。所以这三种运算全都可以用一种简单的谢弗元件表示。

\begin{figure}[h!]
\centering
\includegraphics[width=0.7\textwidth]{图18.5.png} 
\caption{}
\label{fig:18_5}
\end{figure}

\begin{figure}[h!]
\centering
\includegraphics[width=0.7\textwidth]{图18.6.png} 
\caption{}
\label{fig:18_6}
\end{figure}

关于基本元件的选择方法自然并不是唯一的。除了谢弗元件外还可以选用其他的元件。但是，我们选择了谢弗元件，可以使以后的讨论方便。我们将首先考虑用谢弗元件表示的一个运算中的误差控制问题；然后，由谢弗元件组成的任何复杂运算也可以类似地进行设计。



\section{复合方法}

利用增加元件数目来改善可靠性的这一个概念，我们用包含有 $n$ 根输入的一束输入代替一个输入。这样一来，在单独一个谢弗元件的系统中，对于 $a$ 输入线我们将有 $n$ 根线，写作 $a_i$，其中 $i=1,2,\dots,n$；对于 $b$ 输入线也有 $n$ 根，写作 $b_i$，其中 $i=1,2,\dots,n$；输出也是有 $n$ 根的一束线，然后我们规定一个分数 $\delta$，$0<\delta<\dfrac{1}{2}$，如果输出线束中超过 $(1-\delta)n$ 根线是开（或者关），我们就认为从整体来看输出就是开（或者关）。如果输出线束中有不到 $\delta n$ 根线是开（或者关），作为一个整体来看就被考虑为关（或者开）。其他任何中间情况都当作是失灵现象。$\delta$ 就是\underline{置信水平}。问题在于如何利用谢弗元件组成一个系统，当输入线束中误差的概率一定，以及单独元件发生失灵现象的概率一定的时候，可以使得整个系统发生\underline{失灵的概率}减小。

\begin{figure}[h!]
\centering
\includegraphics[width=0.7\textwidth]{图18.7.png} 
\caption{}
\label{fig:18_7}
\end{figure}

作为问题的初步处理，我们从输入线束 $a$ 中取出一根线 $a_i$，从输入线束 $b$ 中取出一根线 $b_i$，把这两根线当作一个谢弗元件的输入线。系统的组成情况像图\ref{fig:18_7}那样。很明显，如果两束输入线中差不多所有的线都是开，那么输出线束中差不多所有的线都是关。如果两束输入线中差不多所有的线都是关，那么输出线束中差不多所有的线都是开。这样一个总的运转状态看起来似乎是合乎需要的，但是更仔细地考虑一下就会发现实际情况并不是像上面所说那样。因为，如果要求谢弗元件的输出是关，那么两根输入线都必须是开才行。在 $a$ 束或者 $b$ 束中发生一个误差将使得输出线束中也产生一个误差。所以，对于输出是关的情形来说在输出线束中产生的误差相当于输入线束中所产生的那些误差的\underline{总和}。同样，在只有一根输入是关而输出是开的情况下，输出线束中所产生的误差和输入线束中产生的误差\underline{相等}。在两根输入线都是关而输出是开的情况下，只有两束输入线束\underline{同时}产生误差那么输出才产生误差；所以，输出线束发生误差的情形比较输入线束发生误差的情况要少。因此，输出的误差和输入的误差并不是一致，有时误差增加，也有时误差减少。于是，就发生了误差的分散情况。这是我们所不希望发生的情况，因为误差的分散现象会使得输出线束中激发线的数目在 $\delta n$ 根线和 $(1-\delta)n$ 根线这个区域中间，这样一来就增加了系统运转失灵的机会。

为了抵制误差的分散现象，我们按照以下的方式引入系统的一个\underline{复原机构}，我们从图\ref{fig:18_7}所表示的\underline{执行机构}中（“执行”的意思就是整个系统的这种用谢弗元件表示的运算方式）取出每一根输出线，然后把它分成两根线，于是我们得到了 $2n$ 根线，然后再把这 $2n$ 根线排列起来，排列的次序是随机的。相继取出一对对的线来，把这样的每一对线作为一个谢弗元件的输入，于是我们又得到包含有 $n$ 根的一束输出线。这种做法用图表示成图\ref{fig:18_8}的样子。如果原来的线束中有 $\alpha_0n$ 根激发线，那么，在统计的意义上输出中未激发线所占的百分数是 $\alpha_0\alpha_0=\alpha_0^2$。输出中激发线所占的百分数就是
\begin{equation}
\alpha_1 = 1-\alpha_0^2。
\label{eq:ch18_single_recovery_unit_alpha1}
\end{equation}

\begin{figure}[h!]
\centering
\includegraphics[width=0.5\textwidth]{图18.8.png} 
\caption{}
\label{fig:18_8}
\end{figure}

假设开初的那些线中激发的概率是 $\alpha_0$，那么，只要 $n$ 很大，使得变换后的激发线的概率就是 $\alpha_1$，所以这还不是一个复原机构。但是，如果我们把这样的两个单元串联起来，那么最后那些线激发的概率 $\alpha_2$ 可以写成
\begin{equation}
\alpha_2 = 1-\alpha_1^2 = 1-(1-\alpha_0^2)^2 = 2\alpha_0^2-\alpha_0^4。
\label{eq:ch18_double_recovery_unit_alpha2}
\end{equation}

\begin{wrapfigure}{r}{0.25\textwidth}
    \centering
    \includegraphics[width=0.25\textwidth]{图18.9.png}
    \caption{ }
    \label{fig:18_9}
\end{wrapfigure}
由于以下的理由，可以知道串联起来的两个单元是一个复原机构，图\ref{fig:18_9}表示出 $\alpha_0$ 和 $\alpha_2$ 之间的关系，当下面关系式成立的时候 $\alpha_2$ 等于 $\alpha_0$，
\begin{equation*}
\alpha_0^4-2\alpha_0^2+\alpha_0 = 0，
\end{equation*}
或者说，$\alpha_0=0,\,\frac{1}{2}(\sqrt{5}-1)$，或者 $1$。所以，如果 $\alpha_0$ 在 $0$ 和 $\frac{1}{2}(\sqrt{5}-1)$ 之间，
\begin{equation*}
\frac{1}{2}(\sqrt{5}-1)=0.618034，
\end{equation*}
$\alpha_2$ 就比 $\alpha_0$ 小；如果 $\alpha_0$ 在 $\frac{1}{2}(\sqrt{5}-1)$ 和 $1$ 之间，$\alpha_2$ 就比 $\alpha_0$ 大，因此复原机构的作用在于使得输出中激发线的概率趋向 $0$ 和 $1$，因此，就减小了执行机构所引起的误差分散现象。

根据上面的讨论的基础，控制误差的系统就是由包含 $n$ 个谢弗元件的一个执行机构和一个复原机构组成的，复原机构是由图\ref{fig:18_7}表示的两个单元（每个单元包括 $n$ 个谢弗元件和一个随机排列的装置）组成。所以，对于每个谢弗元件都完全精确的情况，我们必须把系统扩大到包含 $3n$ 个谢弗元件。我们看到，当给定了某个置信水平 $\delta$，给出输入线束中产生误差的概率以及谢弗元件产生误差的概率以后，用增加元件数目 $n$ 的办法就可以使系统的可靠程度提高到我们所需要的任何程度。所以，基本上，我们还是采用了把元件加多的原理。但是我们在这一节里的分析，给出怎样把这些元件组合起来的办法。冯·诺伊曼把这种用许多不可靠的元件组合成一个非常可靠的系统的特殊方法，叫做“复合方法”。整个系统就称为“复合系统”。

\section{执行机构中的误差}

现在我们来计算由谢弗元件组合起来的复合系统的误差，具体的组合情况上一节已经叙述过了。现在我们来证明误差的确受到了控制。首先，我们看到，在执行机构或者复原机构中，误差的直接来源是每一个谢弗元件本身。假设每个谢弗元件发生错误的概率是 $\varepsilon$。假设有 $r$ 个并联的谢弗元件，并且它们的运转相互无关，因而系统中元件失效的概率仍然是 $\varepsilon$。所以在 $r$ 个元件中，平均有 $\varepsilon r$ 个会发生错误，这个数目也是 $r$ 个并联起来的元件中，最可能发生错误的数目，其他数目的元件发生错误的概率比较小。实际上，由 $r$ 个元件并联而成的系统中，如果每个元件失效的概率等于 $\varepsilon$，那么，决定有 $\rho$ 个元件发生错误的概率 $p_0(\rho,\varepsilon,r)$ 这种问题是随机采样的古典问题。当 $r$ 很大的时候，已经知道\footnote{可以参阅 H. Margenau and G. M. Murphy, \textit{The Mathematics of Physics and Chemistry}, p.422. D. Van Norstrand Company. Inc. , New York, (1943)。}
\begin{equation}
p_0(\rho,\varepsilon,r) \approx \frac{1}{\sqrt{2\pi}\sqrt{\varepsilon(1-\varepsilon)r}} \mathrm{e}^{-\frac{1}{2}\frac{(\rho-\varepsilon r)^2}{\varepsilon(1-\varepsilon)r}}。
\label{eq:ch18_component_error_normal_dist}
\end{equation}
所以概率分布 $p_0(\rho,\varepsilon,r)$ 是按平均值等于 $\varepsilon r$ 而方差等于 $\varepsilon(1-\varepsilon)r$ 的正态分布。

执行机构中，误差的另外一个来源等于从输入线束到各个谢弗元件的那些线配置不当而产生。例如，图\ref{fig:18_7}的输入线束 $a$ 中，如果激发线的百分数等于 $\xi$；输入线束 $b$ 中，如果激发线的百分数等于 $\eta$。我们希望输出线束中有同样百分数的抑制线。但是如果某 $i$ 个元件的输入线 $a_i$ 激发了，而 $b_i$ 未激发，甚至于元件的其他部分并没有毛病，第 $i$ 个元件的输出仍然是激发的情形。用 $\zeta$ 表示执行机构的输出线束中激发线所占的百分数。于是输出线束中实际的抑制线有 $(1-\zeta)n$。$a$ 束中所有激发线的数目等于 $\xi n$，在 $b$ 束中等于 $\eta n$。有效的数目，或者配置适当的数目，$a$ 束中只是 $(1-\zeta)n$，其余 $[\xi-(1-\zeta)]n$ 根线没有生效。$b$ 束中没有生效的数目是 $[\eta-(1-\zeta)]n$，所以有效的输出线等于 $(1-\zeta)n$，由于 $a$ 束中激发线配置不当而有 $[\xi-(1-\zeta)]n$ 根输出线没有生效，由于 $b$ 束中激发线配置不当而有 $[\eta-(1-\zeta)]n$ 根输出线没有生效，最后由于未激发输入线而使输出线没有生效的数目是
\begin{equation*}
\{1-(1-\zeta)-[\xi-(1-\zeta)]-[\eta-(1-\zeta)]\}n = (2-\xi-\eta-\zeta)n。
\end{equation*}
因此，这一类输出可能的有效组合的数目是\footnote{参看 H. Margenau and G. M. Murphy 的书，第 415 页。}
\begin{equation*}
\frac{n!}{[(1-\zeta)n]!\{[\xi-(1-\zeta)]n\}!\{[\eta-(1-\zeta)]n\}![(2-\xi-\eta-\zeta)n]!}。
\end{equation*}
另一方面，在输入方面，$a$ 输入线束中有 $\xi n$ 根激发线，$(1-\xi)n$ 根未激发线，这 $n$ 根线可能的组合数是
\begin{equation*}
\frac{n!}{(\xi n)![(1-\xi)n]!},
\end{equation*}
$b$ 输入线束中，有 $\eta n$ 根激发线，$(1-\eta)n$ 根未激发线，这 $n$ 根线可能的组合数是
\begin{equation*}
\frac{n!}{(\eta n)![(1-\eta)n]!},
\end{equation*}
如果每一个谢弗元件是不发生毛病的理想元件，两束输入中激发线的百分数分别是 $\xi$ 和 $\eta$，而输出中激发线的百分数等于 $\zeta$ 的概率 $p_1$ 可以由下面公式得到
\begin{equation}
\begin{aligned}
p_1(\xi,\eta,\zeta;n) &= \frac{\displaystyle \frac{n!}{[(1-\zeta)n]!\{[\xi-(1-\zeta)]n\}!\{[\eta-(1-\zeta)]n\}![(2-\xi-\eta-\zeta)n]!}}{\displaystyle \frac{n!}{[\xi n]![(1-\xi)n]!} \cdot \frac{n!}{[\eta n]![(1-\eta)n]!}} \\
&= \frac{[\xi n]![(1-\xi)n]![\eta n]![(1-\eta)n]!}{[(1-\zeta)n]!\{[\xi-(1-\zeta)]n\}!\{[n-(1-\zeta)]n\}![(2-\xi-\eta-\zeta)n]!n!}。
\end{aligned}
\label{eq:ch18_ideal_exec_prob_p1}
\end{equation}
很清楚，为了使得计算有意义，前面一段讨论过的四类输出线不能小于零，只要发生小于零的情况，概率就等于零。也就是说，只要违背了下述条件中的任何一个的时候，$p_1$ 就等于零，
\begin{equation}
\left.
\begin{aligned}
1-\zeta &> 0, \\
\xi-(1-\zeta) &> 0, \\
\eta-(1-\zeta) &> 0, \\
2-\xi-\eta-\zeta &> 0。
\end{aligned}
\right\}
\label{eq:ch18_valid_range_inequalities}
\end{equation}

我们将要在 $n$ 很大的假设下化简方程\eqref{eq:ch18_ideal_exec_prob_p1}。当 $n$ 很大，由斯特灵(Stirling)公式，阶乘($!$)可以用它的渐近值表示
\begin{equation}
n! \approx \sqrt{2\pi} \mathrm{e}^{-n} n^{n+\frac{1}{2}}。
\label{eq:ch18_stirling_approx}
\end{equation}
利用方程\eqref{eq:ch18_stirling_approx}我们可以写出 $p_1$ 的近似式
\begin{equation}
p_1(\xi,\eta,\zeta;n) = \frac{1}{\sqrt{2\pi n}} \sqrt{a}\, \mathrm{e}^{-\theta n},
\label{eq:ch18_p1_asymptotic_form}
\end{equation}
其中
\begin{equation}
a = \frac{\xi(1-\xi)\eta(1-\eta)}{(\xi+\zeta-1)(\eta+\zeta-1)(1-\zeta)(2-\xi-\eta-\zeta)}
\label{eq:ch18_def_a_coeff}
\end{equation}
以及
\begin{equation}
\begin{aligned}
\theta &= (\xi+\zeta-1)\log(\xi+\zeta-1) + (\eta+\zeta-1)\log(\eta+\zeta-1) \\
&\quad + (1-\zeta)\log(1-\zeta) + (2-\xi-\eta-\zeta)\log(2-\xi-\eta-\zeta) \\
&\quad - \xi\log\xi - (1-\xi)\log(1-\xi) - \eta\log\eta - (1-\eta)\log(1-\eta)。
\end{aligned}
\label{eq:ch18_def_theta_function}
\end{equation}
把 $\theta$ 对 $\zeta$ 求微商，我们得到
\begin{equation}
\frac{\partial \theta}{\partial \zeta} = \log\frac{(\xi+\zeta-1)(\eta+\zeta-1)}{(1-\zeta)(2-\xi-\eta-\zeta)},
\label{eq:ch18_theta_first_deriv_zeta}
\end{equation}
以及
\begin{equation}
\frac{\partial^2 \theta}{\partial \zeta^2} = \frac{1}{\xi+\zeta-1} + \frac{1}{\eta+\zeta-1} + \frac{1}{1-\zeta} + \frac{1}{2-\xi-\eta-\zeta}。
\label{eq:ch18_theta_second_deriv_zeta}
\end{equation}
由于这些方程我们找到在 $\zeta=1-\xi\eta$，$\theta=\frac{\partial \theta}{\partial \zeta}=0$。更进一步，因为方程\eqref{eq:ch18_valid_range_inequalities}的条件，$\frac{\partial^2 \theta}{\partial \zeta^2}$ 总是正数，所以 $\theta$ 只有一个零点在 $\zeta=1-\xi\eta$。于是假设 $n$ 很大，方程\eqref{eq:ch18_p1_asymptotic_form}中的负幂说明我们只需要考虑 $\theta$ 在它的零点附近的情况就够了。但是，在 $\theta$ 的零点，$\zeta=1-\xi\eta$，
\begin{equation*}
\frac{\partial^2 \theta}{\partial \zeta^2} = \frac{1}{\xi(1-\eta)} + \frac{1}{\eta(1-\xi)} + \frac{1}{\xi\eta} + \frac{1}{(1-\xi)(1-\eta)} = \frac{1}{\xi(1-\xi)\eta(1-\eta)}。
\end{equation*}
在 $\zeta=1-\xi\eta$ 的附近，$\theta$ 可以近似的写成
\begin{equation}
\theta \sim \frac{1}{2}\frac{[\zeta-(1-\xi\eta)]^2}{\xi(1-\xi)\eta(1-\eta)}。
\label{eq:ch18_theta_quadratic_approx}
\end{equation}
$a$ 是 $\zeta$ 的函数，当 $n$ 很大的时候，它和方程\eqref{eq:ch18_p1_asymptotic_form}的指数比较起来是一个变化缓慢的函数。所以，我们可以取 $a$ 在点 $\zeta=1-\xi\eta$ 的值，或者
\begin{equation}
a \sim \frac{1}{\xi(1-\xi)\eta(1-\eta)}。
\label{eq:ch18_a_approx_at_min}
\end{equation}
因此，最后当 $n$ 很大时 $p_1(\xi,\eta,\zeta;n)$ 的近似表达式是
\begin{equation}
p_1(\xi,\eta,\zeta;n) \approx \frac{1}{\sqrt{2\pi\xi(1-\xi)\eta(1-\eta)n}} \mathrm{e}^{-\frac{1}{2}\left[\frac{\zeta-(1-\xi\eta)}{\sqrt{\xi(1-\xi)\eta(1-\eta)/n}}\right]^2}。
\label{eq:ch18_p1_gaussian_approx_final}
\end{equation}
所以 $p_1$ 也是 $\zeta$ 的正态分布函数。

当 $n$ 很大的时候，我们来把概率表达式 $p_1(\xi,\eta,\zeta;n)$ 进行修改，而得到连续的分布函数 $W(\zeta;\xi,\eta;n)$。假设 $W(\zeta;\xi,\eta;n)\mathrm{d}\zeta$ 表示输出中激发线的数目在 $\zeta n$ 与 $\zeta n+1=n(\zeta+1/n)$ 之间的概率，于是 $\mathrm{d}\zeta=1/n$，而这个概率就等于 $p_1(\xi,\eta,\zeta;n)$。所以
\begin{equation}
W(\zeta;\xi,\eta;n) = n p_1 = \frac{1}{\sqrt{2\pi\xi(1-\xi)n(1-\eta)/n}} \mathrm{e}^{-\frac{1}{2}\left[\frac{\zeta-(1-\xi\eta)}{\sqrt{\xi(1-\xi)\eta(1-\eta)/n}}\right]^2}。
\label{eq:ch18_continuous_dist_W}
\end{equation}
于是 $W(\zeta;\xi,\eta;n)\mathrm{d}\zeta$ 通常表示当输入线束中激发线的百分数分别是 $\xi$ 和 $\eta$ 时，输出线束中激发线的百分数在 $\zeta$ 和 $\zeta+\mathrm{d}\zeta$ 之间的概率。线束的大小由线数 $n$ 决定。$W$ 是一个平均值是 $1-\xi\eta$ 而方差等于 $\xi(1-\xi)\eta(1-\eta)/n$ 的高斯分布。可以用等价的方式表达这个结果，写成
\begin{equation}
\zeta = (1-\xi\eta) + \sqrt{\frac{\xi(1-\xi)\eta(1-\eta)}{n}}\,y,
\label{eq:ch18_zeta_gaussian_repr}
\end{equation}
其中 $y$ 表示一个随机变量，这个变量的分布是平均值等于 $0$，而方差等于 $1$ 的标准高斯分布。方程\eqref{eq:ch18_zeta_gaussian_repr}表示 $\zeta$ 的分布是平均值等于 $1-\xi\eta$ 而方差等于 $\xi(1-\xi)\eta(1-\eta)/n$ 的高斯分布。所以方程\eqref{eq:ch18_continuous_dist_W}和方程\eqref{eq:ch18_zeta_gaussian_repr}表示同一个事实，但是方程\eqref{eq:ch18_zeta_gaussian_repr}用起来更加方便些。

我们现在把误差的两种来源合在一起考虑，把各个不完善的谢弗元件的影响加到 $\zeta$ 的分布函数方程\eqref{eq:ch18_zeta_gaussian_repr}上。和方程\eqref{eq:ch18_component_error_normal_dist}相似，方程\eqref{eq:ch18_component_error_normal_dist}可以写成
\begin{equation}
\rho = \varepsilon r + \sqrt{\varepsilon(1-\varepsilon)r}\,y。
\label{eq:ch18_rho_error_gauss}
\end{equation}
在执行机构中有两类谢弗元件。假定输出是激发线的 $n\zeta$ 个元件算一类。发生一个错误就会把激发线的数目减少一个。在这一类元件中 $q$ 个元件发生错误的概率由方程\eqref{eq:ch18_rho_error_gauss}给出
\begin{equation}
q = \varepsilon \zeta n + \sqrt{\varepsilon(1-\varepsilon)\zeta n}\,y。
\label{eq:ch18_q_error_count}
\end{equation}
此外，假定输出是未激发的 $(1-\zeta)n$ 个元件算作另一类。一个错误将会增加一根激发输出线，这一类元件中有 $q'$ 个发生错误的概率分布是
\begin{equation}
q' = \varepsilon(1-\zeta)n + \sqrt{\varepsilon(1-\varepsilon)(1-\zeta)n}\,y,
\label{eq:ch18_q_prime_error_count}
\end{equation}
所以 $q'-q$ 表示由于元件本身的误差使激发输出线增加的数目。根据方程\eqref{eq:ch18_q_error_count}和\eqref{eq:ch18_q_prime_error_count}，我们得到
\begin{equation}
q'-q = 2\varepsilon\left(\frac{1}{2}-\zeta\right)n + \sqrt{\varepsilon(1-\varepsilon)(1-\zeta)n}\,y - \sqrt{\varepsilon(1-\varepsilon)\zeta n}\,y,
\label{eq:ch18_delta_q_expression}
\end{equation}
上面方程中最后两项是两个按照正态分布的随机变量的差。我们可以证明两个按正态分布的随机变量的差仍然是一个按照正态分布的随机变量。

考虑两个按照正态分布的随机变量 $z_1$ 和 $z_2$，它们的平均值都是 $0$，方差分别等于 $\sigma_1^2$ 和 $\sigma_2^2$，因而
\begin{equation}
\left.
\begin{aligned}
z_1 &= \sigma_1 y, \\
z_2 &= \sigma_2 y。
\end{aligned}
\right\}
\label{eq:ch18_two_normal_vars_def}
\end{equation}
或者，用 $W_1(z_1)$ 表示 $z_1$ 的概率分布函数，$W_2(z_2)$ 表示 $z_2$ 的概率分布函数，
\begin{equation*}
W_1(z_1) = \frac{1}{\sigma_1\sqrt{2\pi}} \mathrm{e}^{-\frac{1}{2}\left(\frac{z_1}{\sigma_1}\right)^2},
\end{equation*}
\begin{equation*}
W_2(z_2) = \frac{1}{\sigma_2\sqrt{2\pi}} \mathrm{e}^{-\frac{1}{2}\left(\frac{z_2}{\sigma_2}\right)^2}。
\end{equation*}
现在假设两个随机变量相互独立，$z_1$ 在 $z_1$ 与 $z_1+\mathrm{d}z_1$ 之间，$z_2$ 在 $z_2$ 与 $z_2+\mathrm{d}z_2$ 之间的联合概率是
\begin{equation*}
W_1(z_1)W_2(z_2)\mathrm{d}z_1\mathrm{d}z_2。
\end{equation*}
我们现在引入新的变量 $x_1$ 和 $x_2$，它们定义如下：
\begin{equation*}
\begin{aligned}
x_1 &= z_1 - z_2, \\
x_2 &= z_1 + z_2, \\
z_1 &= \frac{1}{2}(x_1 + x_2), \\
z_2 &= \frac{1}{2}(x_2 - x_1)。
\end{aligned}
\end{equation*}
新变量的联合概率是
\begin{equation*}
\frac{1}{2}W_1\left(\frac{x_1+x_2}{2}\right)W_2\left(\frac{x_2-x_1}{2}\right)\mathrm{d}x_1\mathrm{d}x_2。
\end{equation*}
把这个联合概率从 $-\infty$ 到 $\infty$ 对 $x_2$ 积分，我们就得到概率 $W_1(x_1)\mathrm{d}x_1$，其中 $W_1(x_1)$ 是 $x_1=z_1-z_2$ 的概率分布函数。这样
\begin{equation*}
\begin{aligned}
W(x_1) &= \frac{1}{2}\int_{-\infty}^{\infty} W_1\left(\frac{x_1+x_2}{2}\right)W_2\left(\frac{x_2-x_1}{2}\right)\mathrm{d}x_2 \\
&= \frac{1}{4\pi\sigma_1\sigma_2}\int_{-\infty}^{\infty} \mathrm{e}^{-\frac{1}{2}\left[\left(\frac{x_1+x_2}{2\sigma_1}\right)^2+\left(\frac{x_2-x_1}{2\sigma_2}\right)^2\right]} \mathrm{d}x_2 \\
&= \frac{1}{4\pi\sigma_1\sigma_2}\mathrm{e}^{-\frac{1}{2}\frac{x_1^2}{\sigma_1^2+\sigma_2^2}} \int_{-\infty}^{\infty} \mathrm{e}^{-\xi^2} \frac{\mathrm{d}\xi}{\sqrt{(1/8\sigma_1^2)+(1/8\sigma_2^2)}} \\
&= \frac{1}{\sqrt{2\pi}} \frac{1}{\sqrt{\sigma_1^2+\sigma_2^2}}\mathrm{e}^{-\frac{1}{2}\left(\frac{x_1}{\sqrt{\sigma_1^2+\sigma_2^2}}\right)^2}。
\end{aligned}
\end{equation*}
所以我们可以写
\begin{equation}
z_1 - z_2 = \sqrt{\sigma_1^2+\sigma_2^2}\,y。
\label{eq:ch18_normal_var_diff_rule}
\end{equation}
把联合概率对 $x_1$ 求积分，我们可以得到
\begin{equation}
z_1 + z_2 = \sqrt{\sigma_1^2+\sigma_2^2}\,y,
\label{eq:ch18_normal_var_sum_rule}
\end{equation}
因此两个按正态分布的相互独立随机变量，它们的和以及差仍然是按正态分布的随机变量，新概率分布函数的方差等于原来两个随机变量的方差的和。按正态分布而且平均值等于 $0$ 的随机变量相加或者相减所具有的这种特性是可以预料到的，因为它们对于正值和负值概率相等。

借助于方程\eqref{eq:ch18_normal_var_diff_rule}的关系式，我们可以写出方程\eqref{eq:ch18_delta_q_expression}的结果
\begin{equation*}
q'-q = 2\varepsilon\left(\frac{1}{2}-\zeta\right)n + \sqrt{\varepsilon(1-\varepsilon)n}\,y。
\end{equation*}
令 $(q'-q)/n=\Delta\zeta$ 表示对于不完善的谢弗元件所占的百分数 $\zeta$ 的修正量，我们就得到
\begin{equation}
\Delta\zeta = 2\varepsilon\left(\frac{1}{2}-\zeta\right) + \sqrt{\frac{\varepsilon(1-\varepsilon)}{n}}\,y,
\label{eq:ch18_zeta_correction_delta}
\end{equation}

现在我们可以把方程\eqref{eq:ch18_zeta_gaussian_repr}与\eqref{eq:ch18_zeta_correction_delta}结合起来，于是修正了的输出线中激发线的百分数 $\zeta'$ 是
\begin{equation}
\begin{aligned}
\zeta' &= \zeta + \Delta\zeta = \zeta + 2\varepsilon\left(\frac{1}{2}-\zeta\right) + \sqrt{\frac{\varepsilon(1-\varepsilon)}{n}}\,y \\
&= (1-\xi\eta) + 2\varepsilon\left(\xi\eta - \frac{1}{2}\right) + (1-2\varepsilon)\sqrt{\frac{\xi(1-\xi)\eta(1-\eta)}{n}}\,y \\
&\quad + \sqrt{\frac{\varepsilon(1-\varepsilon)}{n}}\,y,
\end{aligned}
\label{eq:ch18_zeta_prime_combined_18_25}
\end{equation}
方程\eqref{eq:ch18_zeta_prime_combined_18_25}的最后两项表示两个按正态分布的随机变量相加。这样，我们也就可以用方程\eqref{eq:ch18_normal_var_sum_rule}了。所以，最后，用 $\zeta$ 代替 $\zeta'$，从方程\eqref{eq:ch18_zeta_correction_delta}我们得到
\begin{equation}
\zeta = (1-\xi\eta) + 2\varepsilon\left(\xi\eta - \frac{1}{2}\right) + \sqrt{\frac{(1-2\varepsilon)^2\xi(1-\xi)\eta(1-\eta) + \varepsilon(1-\varepsilon)}{n}}\,y,
\label{eq:ch18_final_zeta_distribution_18_26}
\end{equation}
其中 $y$ 是按正态分布的随机变量，平均值等于 $0$ 而方差等于 $1$。方程\eqref{eq:ch18_final_zeta_distribution_18_26}描述了由谢弗元件组成的复合系统中执行机构的性质，这个系统中输入激发线的百分数分别是 $\xi$ 和 $\eta$，输出激发线的百分数是 $\zeta$，各个谢弗元件失效的概率等于 $\varepsilon$。


\section{复合系统的误差}

分析了由谢弗元件组成的复合系统中执行机构的性质以后，我们进行其他的计算就非常容易了。复原机构的每一个单元（图\ref{fig:18_8}）实际上等价于一个执行机构。对于复原机构的第一个单元，输入线是把执行机构的输出线分开而成的。所以我们可以只用同一个百分数 $\xi$ 来代替两个不同的百分数 $\xi$ 和 $\eta$，因此，假设 $\mu$ 是第一个单元输出激发线的百分数，于是根据方程\eqref{eq:ch18_final_zeta_distribution}
\begin{equation}
\mu = (1-\zeta^2) + 2\varepsilon\left(\zeta^2 - \frac{1}{2}\right) + \sqrt{\frac{(1-2\varepsilon)^2\zeta^2(1-\zeta^2)+\varepsilon(1-\varepsilon)}{n}}\,y。
\label{eq:ch18_recovery_unit1_mu_dist}
\end{equation}

与此类似，假设 $v$ 是复原机构中第二个单元的输出中激发线的百分数，那么
\begin{equation}
v = (1-\mu^2) + 2\varepsilon\left(\mu^2 - \frac{1}{2}\right) + \sqrt{\frac{(1-2\varepsilon)^2\mu^2(1-\mu)^2+\varepsilon(1-\varepsilon)}{n}}\,y。
\label{eq:ch18_recovery_unit2_v_dist}
\end{equation}
方程\eqref{eq:ch18_recovery_unit1_mu_dist}和\eqref{eq:ch18_recovery_unit2_v_dist}中，第一项的形式和方程\eqref{eq:ch18_single_recovery_unit_alpha1}完全一样。附加项是由于不完善的元件以及从误差的统计分布所产生的。

只要给定了 $\xi,\eta,\zeta,\varepsilon$ 和 $n$ 时，我们就可以利用方程\eqref{eq:ch18_final_zeta_distribution}到\eqref{eq:ch18_recovery_unit2_v_dist}来计算 $v$ 的分布函数，$v$ 是谢弗元件组成的整个系统中输出激发线的百分数。把表达式回复到原先的概率分布函数形式。我们可以看得比较清楚，例如，方程\eqref{eq:ch18_final_zeta_distribution}等价于
\begin{equation*}
W(\zeta;\xi,\eta;n) = \frac{\exp\left\{-\frac{1}{2}\left[\frac{\zeta-\left[(1-\xi\eta)+2\varepsilon\left(\xi\eta-\frac{1}{2}\right)\right]}{\sqrt{\frac{(1-2\varepsilon)^2\xi(1-\xi)\eta(1-\eta)+\varepsilon(1-\varepsilon)}{n}}}\right]^2\right\}}{\sqrt{2\pi \frac{(1-2\varepsilon)^2\xi(1-\xi)\eta(1-\eta)+\varepsilon(1-\varepsilon)}{n}}}。
\end{equation*}
只要把 $\zeta,\mu$ 和 $v$ 的联合概率对 $\zeta$ 和 $\mu$ 求积分，于是我们就可以得到 $v$ 的概率分布函数 $W(v;\xi,\eta;n)$，也就是
\begin{equation}
\begin{aligned}
W(v;\xi,\eta;n) &= \frac{1}{(2\pi)^{\frac{3}{2}} \sqrt{\frac{(1-2\varepsilon)^2\xi(1-\xi)\eta(1-\eta)+\varepsilon(1-\varepsilon)}{n}}} \int_{-\infty}^{\infty} \mathrm{d}\mu \\
&\quad \cdot \int_{-\infty}^{\infty} \frac{\mathrm{d}\zeta}{\sqrt{\frac{(1-2\varepsilon)^2\zeta^2(1-\zeta)^2+\varepsilon(1-\varepsilon)}{n}} \sqrt{\frac{(1-2\varepsilon)^2\mu^2(1-\mu)^2+\varepsilon(1-\varepsilon)}{n}}} \\
&\quad \cdot \exp\left\{
-\frac{1}{2}\left[\frac{\zeta-\left\{(1-\xi\eta)+2\varepsilon\left(\xi\eta-\frac{1}{2}\right)\right\}}{\sqrt{\frac{(1-2\varepsilon)^2\xi(1-\xi)\eta(1-\eta)+\varepsilon(1-\varepsilon)}{n}}}\right]^2
\right. \\
&\quad -\frac{1}{2}\left[\frac{\mu-\left\{(1-\zeta^2)+2\varepsilon\left(\zeta^2-\frac{1}{2}\right)\right\}}{\sqrt{\frac{(1-2\varepsilon)^2\zeta^2(1-\zeta)^2+\varepsilon(1-\varepsilon)}{n}}}\right]^2 \\
&\quad \left. -\frac{1}{2}\left[\frac{v-\left\{(1-\mu^2)+2\varepsilon\left(\mu^2-\frac{1}{2}\right)\right\}}{\sqrt{\frac{(1-2\varepsilon)^2\mu^2(1-\mu)^2+\varepsilon(1-\varepsilon)}{n}}}\right]^2
\right\}。
\end{aligned}
\label{eq:ch18_full_system_v_joint_dist}
\end{equation}

我们现在可以证明在适当的条件下，如果 $n$ 增大，那么由谢弗元件组成的系统就能够具有合乎理想的性质。考虑一个给定的置信水平 $\delta$，如果两束输入线束中激发线所占的百分数分别用 $\xi$ 和 $\eta$ 表示，$v$ 表示输出线束中未激发线所占的百分数，要求系统具有理想的性质就必须有：$\xi\geqslant1-\delta$，并且 $\eta\geqslant1-\delta$ 时，有 $v\leqslant\delta$；以及当 $\xi\leqslant\delta$ 并且 $\eta\geqslant1-\delta$ 时，或者当 $\xi\geqslant1-\delta$，并且 $\eta\leqslant\delta$ 时，$v\geqslant1-\delta$；以及 $\xi\leqslant\delta$，并且 $\eta\leqslant\delta$ 时 $v\geqslant1-\delta$。我们假定 $n$ 相当大，$\varepsilon$ 相当小，以至于方程\eqref{eq:ch18_final_zeta_distribution}到\eqref{eq:ch18_recovery_unit2_v_dist}中数量级与 $\varepsilon$ 和 $1/\sqrt{n}$ 相同的项都可以略去不计。于是就有
\[
\zeta\approx(1-\xi\eta),\quad \mu\approx1-\xi^2,\quad v\approx1-\mu^2。
\]
或者当 $n\gg1,\varepsilon\ll1$ 时，
\begin{equation}
v \approx 1 - (2\xi\eta - \xi^2\eta^2)。
\label{eq:ch18_ideal_v_approx_large_n}
\end{equation}
现在令 $\xi=1-\alpha,\eta=1-\beta$，并且 $\alpha,\beta\leqslant\delta$；这样 $\xi\geqslant1-\delta,\eta\geqslant1-\delta$，于是方程\eqref{eq:ch18_ideal_v_approx_large_n}给出
\[
v \approx 2(\alpha^2+\beta^2)+\cdots。
\]
因此 $v=O(\delta^2)$\footnote{$O(\delta^2)$ 表示一个与 $\delta^2$ 同数量级的变量。}。与此类似，如果 $\xi\leqslant\delta$，而 $\eta\geqslant1-\delta$，或者 $\xi\geqslant1-\delta,\eta\leqslant\delta$，方程\eqref{eq:ch18_ideal_v_approx_large_n}给出 $v=1-O(\delta^2)$。更进一步，如果 $\xi\leqslant\delta,\eta\leqslant\delta$ 方程也给出 $v=1-O(\delta^4)$。所以，只要 $\varepsilon$ 和 $\delta$ 很小，当 $n\to\infty$ 时，由谢弗元件按复合法组成的系统就能真正得到极为可靠的性质。

当 $n$ 很大但并不是无穷大时，因为要计算方程\eqref{eq:ch18_full_system_v_joint_dist}中的积分值，计算起来就非常麻烦。虽然积分的渐近值可以用古典的方法决定，但是在这里我们不来进行这种计算了。我们将引用冯·诺伊曼所举出过的一个例子：如果 $\delta=0.07$，也就是线束中至少有百分之九十三的激发线就表示一个“正”的信息；至多有百分之七的激发线就表示一个“负”的信息。这时，他发现，为了控制误差，每一个谢弗元件失灵的概率 $\varepsilon$ 必须比 $0.0107$ 小。如果 $\varepsilon\geqslant0.0107$，整个系统失灵的概率不可能由于 $n$ 的增加而变得任意小，对于 $\varepsilon=0.005$，或者说失效的机会是 $0.5\%$ 的情形，冯·诺伊曼给出了表\ref{table_18_1}的数值结果。从这个表里可以看到，甚至于线束包含了 $1000$ 根线，可靠的性能仍然非常差。实际上，它比原来的 $\varepsilon=1\%$ 还要差，但是把 $n$ 再增加 $25$ 倍以后，就可以得出极为可靠的性能。

\begin{table}
\centering
\caption{}   
\mbox{$\delta=0.07 \epsilon = 0.005$}\\  
\label{table_18_1}
\begin{tabular}{cccc}
\toprule
线的数目 $n$ & 失灵的概率 & 线的数目 $n$ & 失灵的概率 \\
\midrule
1000 & $2.7\times10^{-2}$ & 10000 & $1.6\times10^{-10}$ \\
2000 & $2.6\times10^{-3}$ & 20000 & $2.8\times10^{-19}$ \\
3000 & $2.5\times10^{-4}$ & 25000 & $1.2\times10^{-23}$ \\
5000 & $4\times10^{-6}$ & & \\
\bottomrule
\end{tabular}
\end{table}

对于那些本来由谢弗元件组成的系统，前面讨论过的复合法的技巧也还可以照样应用。我们用 $3n$ 个谢弗元件代替原来系统中的每一个谢弗元件。像前面讨论的情形那样，整个系统的误差可以通过系统各个谢弗元件的误差计算出来。实际上，这种计算非常麻烦。可是，为了估计达到规定的可靠性所需要的元件数目 $n$，我们可以认为整个系统等价于一个谢弗元件，而直接用整个反应的结果。这种做法将在下一节里讨论。

\section{一些例子}

为了得到关于所要求线束大小的一个概念，我们考虑有 $2500$ 个真空管的一个计算机，假设每一个真空管平均 $5$ 微秒开动一次。我们要求机器在发生一个错误以前平均工作 $8$ 个小时。在这段时期内，单独一个真空管作用的次数是
\begin{equation*}
\frac{1}{5}\times 8\times 3600\times 10^6 = 5.76\times 10^9。
\end{equation*}
把每个真空管当作一个谢弗元件，考虑每一个真空管是一个独立的单元，于是所要求发生失灵的概率是 $1/(5.76\times10^9)$。但是系统里有 $2500$ 个真空管，$2500$ 个真空管中任何一个发生错误都意味着机器发生错误。所以考虑每个真空管是系统里的一个单元，发生失灵的概率只是前面谈到的 $1/2500$，或者 $1/(2500\times 5.76\times 10^9)=7\times10^{-14}$。于是我们看到，这个发生失灵的概率等于把 $2500$ 个真空管的系统当成一个谢弗元件而得到的概率。在用复合法组成的系统中，这样做法可以大大的简化，关于要求的线数 $n$ 的分析。

如果我们假定置信水平 $\delta$ 以及真空管失效的概率和表 18.1 内所给定的相同，那么，根据这个表，要达到上述的失灵概率需要 $n=14000$。所以，为了使得机器像所要求的那样可靠，需要复合 $14000$ 次。这表示要用有 $3n=3\times14000=42000$ 个真空管的系统去代替机器中每一个真空管。原来 $2500$ 个真空管的机器现在变成 $105\,000\,000$ 个真空管的庞然大物。很显然这是不实际的。

现在讨论另外一个例子，我们考虑关于人类神经系统组织的合理的数量描述。神经系统所包含神经节的数目，通常认为是 $10^{10}$。但是考虑到感受的神经末梢以及其他可能的更小的自主的单元，这个数目当然是太少了，必须还要大几百倍才对。我们就取 $10^{13}$ 为基本元件的数目，神经节在每秒钟里顶多能承受 $200$ 次刺激。但是每秒受到刺激的平均数要少得多，比如说每秒 $10$ 次。我们将要进一步假定我们的神经系统里发生一个错误就相当严重，因而在人的一生中都不应当发生错误。考虑不发生误差的间隔是 $10\,000$ 年，在这段时期内，包括 $10^{13}$ 个元素的系统受到的总刺激数的数目是
\begin{equation*}
10^{13} \times 10\,000 \times 31\,536\,000 \times 10 = 3.2\times 10^{25}。
\end{equation*}
于是失灵的概率应该是
\[
1/(3.2\times 10^{25}) = 3.2\times 10^{-26}。
\]
再假定那些基本的神经元件有表 18.1 叙述的性质；于是从表内推出 $n=28\,000$。

但是需要修正我们的分析：假设人类的神经系统真正复合 $28\,000$ 次，那么在未复合的系统中基本元件的数目并非上面假定的 $10^{13}$；基本元件的数目将要缩减到 $10^{13}/(3\times28\,000)$，因此，失灵的概率应该增加 $3\times28\,000$ 倍，修正后失灵的概率是 $2.7\times10^{-21}$，于是表 18.1 给出 $n=22\,000$。再重复下去，这个值将不会改变了。

这些例子说明，控制误差的复合方法虽然我们可以想像对于神经系统的微小“元件”的情形可以应用，但是在现在工艺情况下，这种方法对于工程系统的应用还是很不实际的。将来一个明显的发展方向是减小元件的体积和功率。从这个观点出发，半导体制成的晶体管比起真空管来有了巨大的改进。因此复合方法将来也许能变为实际可用。另外一个研究的方向是更加深入地分析控制误差的过程。对于谢弗元件，我们的基本系统里，执行机构和复原机构的组织方式毕竟只是一种可能的组织方式。很幸运地，在证明增加可靠性的可能性方面我们已经获得了这一个初步的成功。当然也还可能有另外一种把很多元件组合起来的办法。利用那种方法仍然可以产生同样的可靠性能，而元件的数目比较少。莫尔(E. F. Moore)和申南\footnote{E. F. Moore, C. E. Shannon, \textit{Reliable Circuits Using Less Reliable Relays}, J. Franklin Inst. 262 No. 3, No. 4, 191～208, 281～297(1956)。}(C. E. Shannon)曾经把继电器当作基本元件，采用一种与复合法不相同的组合方法来组成另一个“大”继电器，他们证明：为了得到一个可靠的“大”继电器所需要的继电器的数目可以大大地减少。例如，在同样的条件下用冯·诺伊曼的方法需要的元件数目是 $60\,000$，可是用他们提出的新方法只要 $100$ 个元件就可以得到同样的可靠性。莫尔和申南并且给出达到某个规定的可靠性所至少需要的元件数目。总而言之，自动控制系统中误差的控制问题的研究还仅仅是一个开始。对于控制工程师而言，关于误差的控制问题仍然还没有切实可用的解决方法。
